\subsection{CUDA -- Compute Unified Device Architecture}

CUDA
\begin{itemize}
    \item first release of CUDA 1.0 June, 2007~\cite{nvidia_corporation_cuda_2007}
    \item current release of CUDA 12.8 February, 2025~\cite{nvidia_corporation_cuda_2025-1}
    \item the current CUDA programming guide February, 2025~\cite{nvidia_corporation_cuda_2025}
    \item first CUDA version (3.0) deprecating device emulation mode March, 2010~\cite{nvidia_corporation_cuda_2010}
\end{itemize}

\begin{itemize}
    \item \acrfull{cuda}
    \item NVIDIA's own vendor-specific framework
    \item C-like API with custom extensions \code{<<<...>>>}
    \item only for NVIDIA \glspl{gpu}
    \item offload compute kernels from host (i.e., classically \glspl{cpu}) to one or more devices
    \item many vendor-provided, highly optimized libraries available (e.g., cuBLAS, cuSPARSE, cuFFT, $\dots$)
    \item interfaces for C, C++, FORTRAN, Python, Matlab, Java
\end{itemize}

\begin{advantagebox}
    \begin{itemize}
        \item practically, the go-to standard to target \glspl{gpu}
        \item many resources available
        \item high optimization potential
    \end{itemize}
\end{advantagebox}

\begin{disadvantagebox}
    \begin{itemize}
        \item not standardized
        \item only NVIDIA \glspl{gpu}
        \item closed-source proprietary
    \end{itemize}
\end{disadvantagebox}
